\documentclass[10pt,a4paper]{amsart}
\usepackage[margin=19mm]{geometry}
\usepackage{amsmath,amssymb,mathtools}
\usepackage[scr=rsfs,cal=euler]{mathalfa}
\usepackage{xfrac}
\usepackage[unicode,hidelinks,bookmarks=false]{hyperref}
\newcommand{\ie}{i.e.\ }
\newcommand{\cf}{cf.\ }
\newcommand{\resp}{respectively\ }
\newcommand{\Subsection}{\subsection}
\providecommand{\nobreakdash}{\nobreak\hbox{-}}
\setlength{\emergencystretch}{2em}
\multlinegap0pt
\usepackage{times}
\usepackage{booktabs}
\usepackage{algorithm}
\usepackage{algpseudocode}
\usepackage{url}
\usepackage{parskip}
\setlength{\parskip}{\baselineskip}
\usepackage{etoolbox}
\newtheorem{theorem}{Theorem}
\newtheorem{proposition}[theorem]{Proposition}
\newtheorem{remark}[theorem]{Remark}
\DeclareMathOperator{\diag}{diag}
\DeclareMathOperator{\rank}{rank}
\DeclareMathOperator*{\argmin}{arg\,min}
\DeclareMathOperator*{\argmax}{arg\,max}
\newcommand{\R}{\mathbb{R}}
\newcommand{\sgm}{\sigma}
\newcommand{\proj}{\mathcal{P}}
\newcommand{\spec}{\mathcal{S}}
\newcommand{\Kset}{\mathcal{K}}
\newcommand{\Cset}{\mathcal{C}}
\newcommand{\Ellipsoid}{\mathcal{E}}
\begin{document}
\title[On feasibility problems with spectral constraints]{On feasibility problems with spectral constraints}
\author{Shravan Mohan}
\date{}
\maketitle
\noindent\emph{Source and licence.} On feasibility problems with spectral constraints. Version arXiv:2607.09254v1 (CC BY 4.0). This material is distributed under the Creative Commons Attribution 4.0 International licence, http://creativecommons.org/licenses/by/4.0/, as recorded in the arXiv OAI licence metadata for this exact version and shown on the abstract page https://arxiv.org/abs/2607.09254v1. Full rights notice: licenses/arxiv-2607.09254-CC-BY-4.0.txt.\par\smallskip
\noindent\textit{Editorial scope.} Counted unit: the original \texttt{theorem} environment \texttt{thm:proj} at source lines 277--288 together with its complete proof at lines 290--327; the delivered document keeps source lines 220--328 so that every referenced label and every citation resolves inside it. Native editable source is the paper's own concatenated TeX; the AMS wrapper is editorial formatting only. No publisher class, upstream script or unrelated artwork is redistributed.
\par\smallskip\noindent\textit{Reference note.} This article ships no compiled bibliography (biblatex); the entries delivered here are composed from the author's own BibTeX (.bib) fields, with no field invented and no entry dropped.
\section{Problem Setup}\label{sec:setup}

\subsection{Notation}
For $X \in \R^{m\times n}$ with $r = \min(m,n)$, let
$X = U \diag(\sgm) V^\top$ denote its thin SVD, with the singular
values sorted as
$\sgm_1 \ge \cdots \ge \sgm_r \ge 0$.
We work throughout in the Frobenius inner product
$\langle X, Y\rangle = \mathrm{tr}(X^\top Y)$
with associated norm $\|\cdot\|_F$.
For a closed convex set $\mathcal{S}$, $\proj_{\mathcal{S}}$ denotes
Euclidean projection.

\subsection{Spectral sets and the projection identity}
We say $\spec \subset \R^r$ is \emph{admissible} if it is closed, convex,
contained in the ordered nonnegative cone
$\R^r_{\ge 0,\downarrow} = \{s : s_1 \ge \cdots \ge s_r \ge 0\}$,
and \emph{permutation invariant on the nonnegative orthant} in the sense
that $s \in \spec \iff Ps \in \spec$ for any permutation $P$ acting on
$\R^r_{\ge 0}$ (and projecting back to the ordered cone via sorting).
Let
\begin{equation}\label{eq:Cset}
\Cset \;=\; \{ X \in \R^{m\times n} : \sgm(X) \in \spec \}.
\end{equation}
By construction $\Cset$ is \emph{orthogonally invariant}: if $X \in
\Cset$ then $QXR \in \Cset$ for any orthogonal $Q \in \R^{m\times m}$
and $R \in \R^{n\times n}$.

\subsection{Reducing the matrix projection to a vector projection}
\label{sec:proof}
The Euclidean projection $\proj_\Cset(A)$ is by definition the
minimiser of
\begin{equation}\label{eq:proj_problem}
\min_{X \in \R^{m\times n}} \; \tfrac{1}{2}\|X - A\|_F^2
\quad \text{s.t.}\quad \sgm(X) \in \spec.
\end{equation}
We now reduce \eqref{eq:proj_problem} to a vector problem and exhibit
an optimal $X^\star$ that shares its singular vectors with $A$. The
key tool is the following classical inequality.

\begin{theorem}[von Neumann's trace inequality~\cite{vonNeumann1937}]
\label{thm:vN}
For any $X, A \in \R^{m\times n}$ with singular values
$\sgm(X)$ and $\sgm(A)$ in $\R^r_{\ge 0,\downarrow}$,
\begin{equation}\label{eq:vN}
\langle X, A\rangle \;=\; \mathrm{tr}(X^\top A)
\;\le\; \sum_{i=1}^{r} \sgm_i(X)\,\sgm_i(A).
\end{equation}
Moreover, equality is attained if and only if $X$ and $A$ admit a
\emph{simultaneous SVD}, i.e.\ there exist orthogonal matrices
$U \in \R^{m\times m}$, $V \in \R^{n\times n}$ such that
\begin{equation}\label{eq:sameUV}
X = U \diag(\sgm(X))\,V^\top, \qquad
A = U \diag(\sgm(A))\,V^\top.
\end{equation}
\end{theorem}


\begin{theorem}[Spectral projection identity~\cite{lewis1995,beck}]
\label{thm:proj}
Let $\spec$ be admissible and let $A \in \R^{m\times n}$ have thin SVD
$A = U \diag(s) V^\top$ with $s = \sgm(A) \in \R^r_{\ge 0,\downarrow}$.
Then problem \eqref{eq:proj_problem} attains its minimum at
\begin{equation}\label{eq:projC}
X^\star \;=\; U\,\diag\!\bigl(\proj_\spec(s)\bigr)\,V^\top,
\end{equation}
where $\proj_\spec$ denotes Euclidean projection on $\R^r$. Equivalently,
$\proj_\Cset(A)$ shares the singular vectors $(U, V)$ of $A$, and its
singular values are obtained by projecting $\sgm(A)$ onto $\spec$.
\end{theorem}

\begin{proof}
Expanding the objective of \eqref{eq:proj_problem} gives
\begin{equation}\label{eq:expand}
\tfrac{1}{2}\|X - A\|_F^2 \;=\; \tfrac{1}{2}\|X\|_F^2 + \tfrac{1}{2}\|A\|_F^2 - \langle X, A\rangle.
\end{equation}
The Frobenius norm is unitarily invariant, so
$\|X\|_F^2 = \sum_i \sgm_i(X)^2$ and similarly for $A$. By the
trace inequality \eqref{eq:vN},
\begin{equation}\label{eq:lb}
\langle X, A\rangle \;\le\; \sgm(X)^\top \sgm(A),
\end{equation}
and combining \eqref{eq:expand} and \eqref{eq:lb},
\begin{align}
\tfrac{1}{2}\|X - A\|_F^2
&\;\ge\;
\tfrac{1}{2}\|\sgm(X)\|_2^2 + \tfrac{1}{2}\|\sgm(A)\|_2^2 - \sgm(X)^\top \sgm(A) \notag \\
&\;=\;
\tfrac{1}{2}\|\sgm(X) - \sgm(A)\|_2^2.
\label{eq:lb2}
\end{align}
For any feasible $X \in \Cset$, the right-hand side of \eqref{eq:lb2}
depends only on $\sgm(X) \in \spec$. Hence the optimal value of
\eqref{eq:proj_problem} is at least
\begin{equation}\label{eq:vec_lower}
\min_{t \in \spec}\;\tfrac{1}{2}\|t - \sgm(A)\|_2^2
\;=\;\tfrac{1}{2}\|\proj_\spec(s) - s\|_2^2.
\end{equation}

We now show that this lower bound is attained by the choice
\eqref{eq:projC}. Let $t^\star = \proj_\spec(s) \in \spec$, and set
$X^\star = U \diag(t^\star) V^\top$ with the same orthogonal matrices
$U, V$ as in the SVD of $A$. Then $\sgm(X^\star) = t^\star \in \spec$,
so $X^\star \in \Cset$. Moreover $X^\star$ and $A$ trivially admit a
simultaneous SVD \eqref{eq:sameUV}, so equality holds in
\eqref{eq:vN}, hence in \eqref{eq:lb}, hence in \eqref{eq:lb2}. The
objective at $X^\star$ therefore equals
$\tfrac{1}{2}\|t^\star - s\|_2^2$, which matches \eqref{eq:vec_lower}.
\end{proof}


\begin{thebibliography}{99}
\bibitem[]{beck} Beck, A.. First-Order Methods in Optimization. (2017)
\bibitem[]{lewis1995} Lewis, A. S.. The convex analysis of unitarily invariant matrix functions. Journal of Convex Analysis 2 (1995) 173--183
\bibitem[]{vonNeumann1937} von Neumann, J.. Some matrix inequalities and metrization of matric-space. Tomsk Univ. Rev. 1 (1937) 286--300
\end{thebibliography}
\end{document}
